\begin{frame}{표를 읽는 법: ``재고'' vs ``유량'', MLP 비교}
  \begin{enumerate}
    \item \kb{파라미터는 뒤로 갈수록 커지지만}(40$\to$296$\to$1,168)
      \kb{FLOPs 최댓값은 conv 2}(165,888) --- FLOPs는 $n_{\text{out}}^2$ 에
      비례하고 conv 2가 공간적으로 가장 크다(24$\times$24).
      conv 3은 채널 2배지만 공간 1/4(10$\times$10)라 오히려 적다
    \item \kb{풀링은 파라미터 0}, 비용도 곱셈-덧셈이 아니라
      \textbf{비교}(4개 중 최댓값) --- FLOPs 합계에서 무시할 만큼 작음
  \end{enumerate}
  \vskip0.5em
  세 합성곱 층을 \textbf{같은 출력 체적}(10$\times$10$\times$16 $=1{,}600$)
  의 \textbf{완전연결층 하나}(784 입력 $\to$ 1,600 뉴런)로 대체하면:
  파라미터 $784\times1{,}600+1{,}600 = 1{,}256{,}000$
  (합성곱 스택의 \textbf{약 835배}), FLOPs $1{,}254{,}400$
  (\textbf{약 4.1배}).
  \vskip0.4em
  \begin{block}{``재고'' vs ``유량''}
    \textbf{파라미터는 ``재고''}(층당 고정 저장),
    \textbf{FLOPs는 ``유량''}(출력 위치마다 반복되는 계산) ---
    ``한 번에'' 보는 대신 \textbf{작은 패치를 재사용}하는 합성곱이
    \textbf{약 0.12\%의 파라미터, 1/4의 연산}으로 같은 체적의 특징을
    만든다 --- 게다가 \textbf{위치 불변성}이라는 귀납적 편향까지 덤.
    한쪽만 보고 설계하면 안 된다.
  \end{block}
\end{frame}
